\documentclass[10pt,a4paper]{amsart}
\usepackage[margin=19mm]{geometry}
\usepackage{amsmath,amssymb,mathtools}
\usepackage[scr=rsfs,cal=euler]{mathalfa}
\usepackage{xfrac}
\usepackage[unicode,hidelinks,bookmarks=false]{hyperref}
\newcommand{\ie}{i.e.\ }
\newcommand{\cf}{cf.\ }
\newcommand{\resp}{respectively\ }
\newcommand{\Subsection}{\subsection}
\providecommand{\nobreakdash}{\nobreak\hbox{-}}
\setlength{\emergencystretch}{2em}
\multlinegap0pt
\usepackage{amssymb}
\usepackage{mathtools}
\newtheorem{assumption}{Assumption}
\newtheorem{lemma}{Lemma}
\newcommand{\Hs}{H}
\newcommand{\PP}{\mathcal P}
\newcommand{\uL}{\underline\Lambda}
\newcommand{\un}{\underline\nu}
\title{\bfseries Stabilization of a Heterodirectional\\ Bradytachic (1D, 2D) PDE Pair}
\author{Miroslav Krstic}
\date{Source: arXiv:2609.14941v1 (CC BY 4.0)}
\begin{document}
\maketitle

\noindent\emph{Source and licence.} \bfseries Stabilization of a Heterodirectional\\ Bradytachic (1D, 2D) PDE Pair. Version arXiv:2609.14941v1 (CC BY 4.0), distributed by arXiv under the Creative Commons Attribution 4.0 International licence, http://creativecommons.org/licenses/by/4.0/, as recorded in the arXiv licence metadata for this exact version and shown on the abstract page https://arxiv.org/abs/2609.14941v1. Full licence notice: \texttt{licenses/arxiv-2609.14941-CC-BY-4.0.txt}.\par\smallskip

\noindent\textit{Editorial scope.} Counted unit: the article's lemma lem:evolution (source0.tex lines 727--729). The article's complete original proof of it is delivered with it (source0.tex lines 731--802, one \texttt{proof} environment of 72 lines). No mathematics is written, repaired or re-derived: the statement and the proof are the article's own lines, in the article's own notation, and no retained line is substituted or re-flowed.  Retained uncounted as the source-local context the proof needs: lines 103--108; lines 113--115; lines 146--149; lines 173--176; lines 178--183; lines 193--198; lines 720--725.  Nothing was omitted from any retained span, so the package carries no omission record.  No retained line contains a citation, so the wrapper carries no bibliography and no bibliography entry.  No retained line contains a figure environment, an image-inclusion command or drawing code, so no image is delivered and the package declares no asset dependency.  Editorial formatting only, in addition: an AMS \texttt{amsart} preamble that restates the article's own declarations and macros used by the retained lines, namely the environments \texttt{assumption}, \texttt{lemma} on separate counters, together with the standard packages those lines need.  The retained environments print in this document as Assumption 1, Lemma 1 in order of appearance, whereas the article prints the same items with the numbering of its own full document.  The complete native source of the article as distributed in the arXiv e-print for this exact version is bundled unchanged as \texttt{sources/210374/source0.tex}.
\begin{equation}
 \nu(x,0)\,v(x,0,t)-\varepsilon v_z(x,0,t)=0,
 \qquad
 v_z(x,D,t)=0\quad(\varepsilon>0),
 \label{eq:gen-zbc}
\end{equation}

\begin{assumption}\label{as:coeff}
$\Lambda\in C^1([0,1])$ and $\nu\in C^1([0,1]\times[0,D])$, with $\Lambda(x)\ge\uL>0$ and $\nu(x,z)\ge\un>0$; $b\in C^1([0,1]\times[0,D])$, with $b(x,\cdot)\in C^2([0,D])$ uniformly in $x$; and $c$, $d$, $\theta$ are continuous.
\end{assumption}

\begin{equation}
 V_x(x)=\mathcal A_\varepsilon(x)V(x)+B(x)u(x),\qquad V(0)=0,
 \label{eq:Hspatial}
\end{equation}

\begin{equation}
 (\PP_\varepsilon u)(x)=\int_0^x\Phi_\varepsilon(x,\xi)B(\xi)u(\xi)\,d\xi ,
 \label{eq:Pgen}
\end{equation}

\begin{equation}
 \|\Phi_\varepsilon(x,\xi)\|_{\mathcal L(\Hs)}\le M_\Phi:=e^{\rho_1/2+M_R},
 \qquad
 \|\PP_\varepsilon\|_{\mathcal L(L^2(0,1),L^2_x(\Hs))}\le M_P:=\frac{M_\Phi\sqrt D\,B_0}{\uL} .
 \label{eq:Phibound}
\end{equation}

\begin{equation}
 f_\varepsilon(x,\xi):=C(x)\Phi_\varepsilon(x,\xi)B(\xi),
 \qquad
 |f_\varepsilon(x,\xi)|\le\bar f:=M_CM_P,
 \label{eq:Kgen}
\end{equation}

\begin{equation}
 \partial_z Z(s;x,z)
 =\exp\!\Big(-\!\int_s^x\rho_z(\tau,Z(\tau;x,z))\,d\tau\Big)
 \in\big[e^{-\rho_1(x-s)},\,e^{\rho_1(x-s)}\big].
 \label{eq:charJac}
\end{equation}

\begin{lemma}[Spatial evolution family; continuity of $f_\varepsilon$]\label{lem:evolution}
Under Assumption~\ref{as:coeff}, for every $\varepsilon\in[0,\bar\varepsilon]$ the homogeneous problem $V_x=\mathcal A_\varepsilon(x)V$, $V(\xi)=\phi$, is well posed on $\xi\le x\le1$ for every $\phi\in\Hs$, and its solution operators $\Phi_\varepsilon(x,\xi)\in\mathcal L(\Hs)$ form an evolution family, with $\Phi_\varepsilon(x,x)=I$, $\Phi_\varepsilon(x,s)\Phi_\varepsilon(s,\xi)=\Phi_\varepsilon(x,\xi)$, and $\|\Phi_\varepsilon(x,\xi)\|_{\mathcal L(\Hs)}\le M_\Phi$, uniformly in $\varepsilon$. Moreover $\PP_\varepsilon$ from \eqref{eq:Pgen} is the unique mild solution map of \eqref{eq:Hspatial}, $\|\PP_\varepsilon\|\le M_P$ as an operator from $L^2(0,1)$ to $L^2_x(\Hs)$, uniformly in $\varepsilon$; $\PP_\varepsilon u$ satisfies the homogeneous conditions \eqref{eq:gen-zbc}, and $(\PP_\varepsilon u)(0,z)=0$; and $f_\varepsilon$ from \eqref{eq:Kgen} is continuous on $T$.
\end{lemma}

\begin{proof}
The uniform bound is a dissipativity computation. With $\phi$ in the domain of $\mathcal A_\varepsilon(x)$, for $\varepsilon>0$,
\begin{equation}
 \varepsilon\langle\phi,\phi''\rangle_\Hs
 =\varepsilon\phi(D)\phi'(D)-\varepsilon\phi(0)\phi'(0)-\varepsilon\|\phi'\|_\Hs^2
 =-\nu(x,0)\phi(0)^2-\varepsilon\|\phi'\|_\Hs^2,
 \label{eq:eps-diss1}
\end{equation}
\begin{equation}
 -\langle\phi,\nu\phi'\rangle_\Hs
 =-\tfrac12\nu(x,D)\phi(D)^2+\tfrac12\nu(x,0)\phi(0)^2
 +\tfrac12\langle\nu_z\phi,\phi\rangle_\Hs,
 \label{eq:eps-diss2}
\end{equation}
so the boundary contributions at $z=0$ sum to $-\tfrac12\nu(x,0)\phi(0)^2\le0$ and those at $z=D$ to $-\tfrac12\nu(x,D)\phi(D)^2\le0$; for $\varepsilon=0$, \eqref{eq:eps-diss1} is absent, $\phi(0)=0$, and \eqref{eq:eps-diss2} alone gives the same conclusion. Since $\nu_z/\Lambda=\rho_z$,
\begin{equation}
 \langle\phi,\mathcal A_\varepsilon(x)\phi\rangle_\Hs
 \le\Big(\frac{\rho_1}{2}+M_R\Big)\|\phi\|_\Hs^2,
 \label{eq:eps-diss3}
\end{equation}
and Gronwall applied to $\frac{d}{dx}\|V\|_\Hs\le(\rho_1/2+M_R)\|V\|_\Hs+\|B(x)\|_\Hs|u(x)|$ gives \eqref{eq:Phibound} and the variation-of-constants formula \eqref{eq:Pgen} on classical solutions, extended by density.

Existence of the family: for each fixed $\varepsilon>0$ it follows from Lions' theory of nonautonomous forms. The domain of $\mathcal A_\varepsilon(x)$ varies with $x$ through the Robin coefficient $\nu(x,0)$, but the form domain does not: the associated form on $H^1(0,D)\times H^1(0,D)$,
\begin{equation}
 a_\varepsilon(x;\phi,\psi)
 =\frac{1}{\Lambda(x)}
 \Big[\varepsilon\langle\phi',\psi'\rangle_\Hs
 +\nu(x,0)\phi(0)\psi(0)
 +\langle\nu\phi',\psi\rangle_\Hs
 -\langle d\phi,\psi\rangle_\Hs
 -\langle\Theta\phi,\psi\rangle_\Hs\Big],
 \label{eq:form}
\end{equation}
is bounded and satisfies a G\aa rding inequality on the fixed space $H^1(0,D)$, uniformly in $x$, with continuous dependence on $x$ (measurable and bounded suffices for Lions' construction), so the family exists with solutions in $C([\xi,1];\Hs)$, which gives strong continuity; the $\varepsilon$-uniform bound \eqref{eq:Phibound} then comes from \eqref{eq:eps-diss1}--\eqref{eq:eps-diss3} and not from the generation theorem, whose constants may degenerate as $\varepsilon\downarrow0$. For $\varepsilon=0$ the zero-inflow transport family is explicit,
\begin{equation}
 \big(\Phi_{0,0}(x,\xi)\phi\big)(z)
 =\begin{cases}
 \phi\big(Z(\xi;x,z)\big), & \text{if the characteristic reaches }s=\xi\text{ in }z>0,\\
 0, & \text{if it exits through }z=0\text{ at some }\xi^*\in(\xi,x],
 \end{cases}
 \label{eq:Phi0def}
\end{equation}
with the evolution property from uniqueness of characteristics; substituting $\zeta=Z(\xi;x,z)$ on the survival set and using \eqref{eq:charJac},
\begin{equation}
 \|\Phi_{0,0}(x,\xi)\phi\|_\Hs^2
 =\int\phi(\zeta)^2\,\frac{dz}{d\zeta}\,d\zeta
 \le e^{\rho_1(x-\xi)}\|\phi\|_\Hs^2,
 \label{eq:Phi0bound}
\end{equation}
and the full family at $\varepsilon=0$ is the unique solution of the Volterra operator equation
\begin{equation}
 \Phi_0(x,\xi)=\Phi_{0,0}(x,\xi)+\int_\xi^x\Phi_{0,0}(x,s)\mathcal R(s)\Phi_0(s,\xi)\,ds,
 \label{eq:PhiDuhamel}
\end{equation}
obtained by successive approximation, uniformly convergent in operator norm on $0\le\xi\le x\le1$.

Finally, since $\|B(\xi)\|_\Hs\le\sqrt D\,B_0/\uL$, the kernel $G_\varepsilon(x,\xi):=\Phi_\varepsilon(x,\xi)B(\xi)$ obeys
\begin{equation}
 \|G_\varepsilon(x,\xi)\|_\Hs\le M_\Phi\frac{\sqrt D\,B_0}{\uL}=M_P,
 \label{eq:MG}
\end{equation}
so $\|(\PP_\varepsilon u)(x)\|_\Hs\le M_P\int_0^x|u|\le M_P\|u\|$ and $\|\PP_\varepsilon u\|_{L^2_x(\Hs)}\le M_P\|u\|$. The boundary and initial conditions hold by construction of the mild solution.

Finally, continuity of $f_\varepsilon$: writing
\begin{equation}
 \Phi_\varepsilon(x,\xi)B(\xi)-\Phi_\varepsilon(x',\xi')B(\xi')
 =\Phi_\varepsilon(x,\xi)\big[B(\xi)-B(\xi')\big]
 +\big[\Phi_\varepsilon(x,\xi)-\Phi_\varepsilon(x',\xi')\big]B(\xi'),
 \label{eq:fcont}
\end{equation}
the first term is $O(\|B(\xi)-B(\xi')\|_\Hs)$ by \eqref{eq:Phibound}, with $\xi\mapsto B(\xi)$ continuous into $\Hs$ since $b\in C^1$, and the second tends to zero because the family, applied to the fixed vector $B(\xi')$, is jointly strongly continuous---for $\varepsilon=0$ this is $L^2$-continuity of composition with the characteristic flow \eqref{eq:Phi0def}, unaffected by the moving cut-off, and for $\varepsilon>0$ it is the $C([\xi,1];\Hs)$ regularity above. Pairing \eqref{eq:fcont} with $c(x,\cdot)$, continuous into $\Hs$, gives continuity of $f_\varepsilon$ on $T$; the jump of the $\Hs$-valued kernel in $z$ is integrated out by $C(x)$.
\end{proof}

\end{document}
